\documentclass{article}
\usepackage{amsmath, amssymb, amsthm}
\usepackage{hyperref}
\usepackage{geometry}
\geometry{margin=1in}

\title{Erd\H{o}s Problem \#726}
\date{Last updated: 2025-08-31}
\author{Source: erdosproblems.com}

\begin{document}
\maketitle

\noindent\textbf{Status:} OPEN \quad \textbf{No prize}

\noindent\textbf{Tags:} number theory

\medskip

\noindent\textbf{Problem Statement:}

As $n\to \infty$ ranges over integers\[\sum_{p\leq n}1_{n\in (p/2,p)\pmod{p}}\frac{1}{p}\sim \frac{\log\log n}{2}.\]




A conjecture of Erd\H{o}s, Graham, Ruzsa, and Straus \cite{EGRS75}. For comparison the classical estimate of Mertens states that\[\sum_{p\leq n}\frac{1}{p}\sim \log\log n.\]By $n\in (p/2,p)\pmod{p}$ we mean $n\equiv r\pmod{p}$ for some integer $r$ with $p/2&#60;r&#60;p$.




References

[EGRS75] Erd\H{o}s, P. and Graham, R. L. and Ruzsa, I. Z. and Straus, E. G., On the prime factors of $(\sp{2n}\sb{n})$ . Math. Comp. (1975), 83-92.


\bigskip
\noindent\small{Source: \url{https://www.erdosproblems.com/726}}
\end{document}
